% Source: https://artofproblemsolving.com/community/c3h1928322_questions
% Stable identity: aops-topic-1928322-q01
\textbf{(a)} Since $100=4\cdot25$ with $\gcd(4,25)=1$,
\[
\varphi(100)=\varphi(4)\varphi(25)=(4-2)(25-5)=2\cdot20=\boxed{40}.
\]
So Euler's theorem gives $a^{40}\equiv1\pmod{100}$ for every $a$ relatively prime to $10$.

\textbf{(b)} Let $\gcd(a,10)=1$.
\begin{itemize}
\item $a$ is odd, so $a^2\equiv1\pmod4$ (an odd square is $1$ more than a multiple of $8$). Hence $a^{20}=\left(a^2\right)^{10}\equiv1\pmod4$.
\item $\gcd(a,25)=1$, so by Euler's theorem $a^{\varphi(25)}=a^{20}\equiv1\pmod{25}$.
\end{itemize}
Both $4$ and $25$ divide $a^{20}-1$, and they are relatively prime, so $100\mid a^{20}-1$. Hence $\lambda(100)\le20$.

For the reverse inequality look at $a=3$ modulo $100$:
\[
3^2=9,\quad 3^4=81,\quad 3^5=243\equiv43,\quad 3^{10}\equiv43^2=1849\equiv49,\quad 3^{20}\equiv49^2=2401\equiv1 .
\]
The order of $3$ modulo $100$ divides $20$, and none of $3^1,3^2,3^4,3^5,3^{10}$ is $\equiv1$, so the order is exactly $20$. Every $m$ with $3^m\equiv1\pmod{100}$ is a multiple of $20$, so $\lambda(100)\ge20$. Therefore $\boxed{\lambda(100)=20}$, half of $\varphi(100)$: the exponent in Euler's theorem works but need not be the smallest one.

\textbf{(c)} Since $2019=20\cdot100+19$, part (b) gives $3^{2019}\equiv3^{19}\pmod{100}$. Also $3\cdot3^{19}=3^{20}\equiv1$, so $3^{19}$ is the inverse of $3$ modulo $100$. As $3\cdot67=201\equiv1$, we get $3^{2019}\equiv67\pmod{100}$: the last two digits are $\boxed{67}$.
